\subsection{An indexed family of arithmetic theories}
\label{subsec:indexed-arithmetic-theories}

Put $\mathcal B=\mathcal B_{\PA}$ and let
$I:S_{\PA}\to\Split(S_{\PA})$ be the inclusion into the
restriction splitting. The object $E=I(\omega)=(\omega,1_\omega)$
belongs to $\mathcal B$, and $G_{\N}E=\N$.
Choose a $\Sigma_1$ predicate $\operatorname{Ax}_{\PA}(z)$
for the codes of the chosen sentence axioms of PA. Define
\[
 W(e,z)\equiv\exists y\,G(e,z,y),\qquad
 \operatorname{Ax}(e,z)\equiv
 \operatorname{Sent}(z)\land
       \bigl(\operatorname{Ax}_{\PA}(z)\lor W(e,z)\bigr).
\]
The predicate $W$ enumerates all c.e.\ subsets of $\N$ by
program domains. For $e\in\N$, let $\Gamma_e$ be the set of
sentences whose codes satisfy $\operatorname{Ax}(e,z)$ in $\N$.
These are axiom presentations extending PA. Their internal axiom
subobjects $A_{\Gamma_e}\hookrightarrow S$ use the literal
specializations $\operatorname{Ax}(\bar e,z)$.
The two-variable predicate represents a subobject
\[
 i:A\hookrightarrow E\times S
\]
in $\mathcal B$, through the predicate presentation of
Lemma~\ref{lem:syntactic-total-category}.

Write $j:S\hookrightarrow\Fm$ for the sentence inclusion.
Retain the logical-axiom morphism $\lambda:L\to\Fm$ and
the rule morphisms $p_n:R_n\to\Fm^n$, $q_n:R_n\to\Fm$
for $n=1,2$ from Appendix~\ref{app:first-order-derivations}.
On $F_E=E\times\Fm$, use the rule objects
\[
 R^E_0=(E\times L)+A,\qquad
 R^E_n=E\times R_n\quad(n=1,2).
\]
The nullary premise morphism is the unique map to $1$, and its
conclusion morphism is
\[
 q^E_0=[1_E\times\lambda,(1_E\times j)i]:R^E_0\to F_E.
\]
For $n=1,2$, let
$d_n:E\times\Fm^n\to(E\times\Fm)^n$ be the morphism whose
$k$th component is $\langle\pi_E,\pi_k\pi_{\Fm^n}\rangle$.
Put
\[
 p^E_n=d_n(1_E\times p_n),\qquad
 q^E_n=1_E\times q_n.
\]
Every premise and the conclusion of a rule thus have the same
$E$-coordinate. Lemma~\ref{lem:syntactic-derivation-initiality}
gives the initial derivation algebra with conclusion morphism
\[
 c^{\mathrm{all}}:P^{\mathrm{all}}\longrightarrow E\times\Fm.
\]
Let $c:P\to E\times S$ be its pullback along $1_E\times j$,
and let $m:\mathcal T\hookrightarrow E\times S$ be the regular
image of $c$. Theorem~\ref{thm:least-theorem-subobject} identifies
the image of $c^{\mathrm{all}}$ as the least subobject closed
under the displayed rules. Stability of regular images identifies
$\mathcal T$ with the pullback of that image along $1_E\times j$.
The monomorphism $m$ has a partial inverse in $\Split(S_{\PA})$
by Lemma~\ref{lem:syntactic-total-category}. Hence
Lemma~\ref{lem:range-regular-image} identifies its partial identity
with $t=\ran(c)$.

For $e\in\N$, let $\epsilon_e:1\to E$ be the numeral morphism
in $\mathcal B$, and put
$k_e^X=\langle\epsilon_e!,1_X\rangle:X\to E\times X$
for $X\in\mathcal B$.

\begin{proposition}[Fibers of the derivation family]
\label{prop:arithmetic-theory-fibers}
For every $e\in\N$, the pullback of $c^{\mathrm{all}}$ along
$k_e^{\Fm}$ is isomorphic, as a derivation algebra over $\Fm$,
to the initial derivation algebra for $A_{\Gamma_e}$.
There is an induced morphism $\iota_e:P_{\Gamma_e}\to P$,
and the square
\[
\begin{tikzcd}
 P_{\Gamma_e}\arrow[r,"\iota_e"]
       \arrow[d,"c_{\Gamma_e}"'] & P\arrow[d,"c"]\\
 S\arrow[r,"k_e^S"'] & E\times S
\end{tikzcd}
\]
is a pullback in $\mathcal B$. Consequently
$(k_e^S)^*\mathcal T$ is the sentence-theory subobject
generated by $A_{\Gamma_e}$.
\end{proposition}
\begin{proof}
The numeral morphism $\epsilon_e$ has a complement: put
$E_{\ne e}=(\omega,[x\ne\bar e]_{\PA})\in\mathcal B$,
with its inclusion into $E$. The two carrier predicates $x=\bar e$
and $x\ne\bar e$ give an isomorphism
$E\cong1+E_{\ne e}$ in $\mathcal B$. Products and finite
extensivity, supplied by
Lemma~\ref{lem:syntactic-total-category}, therefore give
\[
 E\times\Fm\cong\Fm+(E_{\ne e}\times\Fm),\qquad
 \Xi_e:\mathcal B/(E\times\Fm)
   \xrightarrow{\simeq}\mathcal B/\Fm\times
                 \mathcal B/(E_{\ne e}\times\Fm).
\]
The first component functor is pullback along $k_e^{\Fm}$.

Let $H_E$ be the inference endofunctor for the family rules,
and let $H_{\ne e}$ be formed by restricting those rules to
parameters in $E_{\ne e}$. There is a natural isomorphism
\[
 \Xi_e H_E\cong
      (H_{A_{\Gamma_e}}\times H_{\ne e})\Xi_e.
\]
To verify it, write a slice object as the coproduct of its two
inverse images. For nullary rules,
$B_0(Q)=R^E_0$, and its two components are the corresponding
inverse images under $q^E_0$. For a rule of arity $n=1,2$,
its premise morphism assigns its parameter to every premise.
Consequently, in the pullback defining $B_n(Q)$, all child
conclusions lie in the same component as the rule conclusion.
The parameter-$e$ component is therefore exactly the pullback
which defines $B_n$ for $A_{\Gamma_e}$; the other component
is the corresponding pullback for $H_{\ne e}$.
These identifications commute with slice morphisms by the
uniqueness of the pullback maps. On the parameter-$e$ component,
the nonlogical axiom predicate is
$\operatorname{Ax}(\bar e,z)$, so its rule system is precisely
the one defining $H_{A_{\Gamma_e}}$.

The equivalence thus induces an equivalence of algebra
categories with
$\operatorname{Alg}(H_{A_{\Gamma_e}})\times
 \operatorname{Alg}(H_{\ne e})$.
Let $(P_e,P_{\ne e})$ be the image of the initial family
algebra. For any $H_{A_{\Gamma_e}}$-algebra $Q$, initiality
provides a morphism to $(Q,P_{\ne e})$, and hence a morphism
$P_e\to Q$. Any two such morphisms, paired with
$1_{P_{\ne e}}$, would give two morphisms from the initial
family algebra. They are equal. Thus $P_e$ is initial and is
isomorphic to the initial derivation algebra for
$A_{\Gamma_e}$ from
Theorem~\ref{thm:syntactic-generated-theory}.

Restricting conclusions to sentences gives the displayed
pullback square. Regular images are stable under pullback in
$\mathcal B$, so the pullback of $\mathcal T$ is the image of
$c_{\Gamma_e}$. Appendix~\ref{app:first-order-derivations}
identifies that image with the generated sentence theory.
\end{proof}

Choose an index $a_c$ representing $c$ in the numerical
predicate presentation, and write
\[
 \operatorname{Prov}(e,z)
       \equiv\exists d\,G_{a_c}(d,\langle e,z\rangle).
\]
This is a $\Sigma_1$ predicate for $\mathcal T$.
Proposition~\ref{prop:arithmetic-theory-fibers} identifies
$\operatorname{Prov}(\bar e,z)$ with the theorem predicate
of the chosen presentation $\Gamma_e$.

\begin{theorem}[Uniform independent sentences]
\label{thm:uniform-independent-sentences}
For the family $\Gamma_e$ above, there is a total morphism
$q:E\to S$ in $\mathcal B$ such that, for every $e\in\N$
for which $\Gamma_e$ is consistent, the sentence with code
$(G_{\N}q)(e)$ is neither provable nor refutable from $\Gamma_e$.
\end{theorem}
\begin{proof}
Use the numeral-substitution morphism $\theta:E\to S$ of
Lemma~\ref{lem:arithmetic-diagonal-representation} and its
numerical function $\Theta$. Put
\[
 C(e,x)=\operatorname{Prov}(e,\Theta(x)),\qquad
 C'(e,x)=\operatorname{Prov}(e,\operatorname{Neg}(\Theta(x))).
\]
Apply Lemma~\ref{lem:syntactic-functional-search} to the
$\Sigma_1$ relation
\[
 R(e,x,z)\equiv
   (C(e,x)\land z=2x+1)\lor(C'(e,x)\land z=2x),
\]
using $\langle e,x\rangle$ as its input code. It gives a single
morphism $\ell:\omega\times\omega\rightharpoonup\omega+\omega$
in $S_{\PA}$ whose graph implies $R$ and whose domain predicate
is $C\lor C'$, with both assertions provable in PA.
On each interpretation in which $C(e,-)$ and $C'(e,-)$ are
disjoint, the graph at parameter $e$ is exactly $R(e,-,-)$,
by the argument of Lemma~\ref{lem:arithmetic-proof-search}.

The parameterized Turing property gives a total morphism
$g:\omega\to\omega$ in $S_{\PA}$ satisfying
\[
 \tau_{\PA}(g\times1_\omega)=b^{-1}\ell.
\]
Precomposition with $\langle1_\omega,g\rangle$ gives
\[
 hg=\ell\langle1_\omega,g\rangle.
\]
Set $q=\theta I(g):E\to S$, which is total.

Fix $e\in\N$ with $\Gamma_e$ consistent, and put
$n=F_{\N}(g)(e)$. Choose $K\models\Gamma_e$ by completeness.
Proposition~\ref{prop:standard-code-soundness} and
Proposition~\ref{prop:arithmetic-theory-fibers} show that
$C(e,-)$ and $C'(e,-)$ have disjoint truth sets in $\N$.
If $C(e,n)$ holds, the interpreted graph of $\ell$ sends
$(e,n)$ to the second copy of $n$. The preceding identity then
puts $n$ in $D_1$. Lemma~\ref{lem:arithmetic-diagonal-representation}
gives $C'(e,n)$, a contradiction. If $C'(e,n)$ holds, the same
identity puts $n$ in $D_0$, and that lemma gives $C(e,n)$,
again a contradiction. Therefore neither predicate holds.
Since $(G_{\N}q)(e)=\Theta(n)$, this is the claimed independence.
\end{proof}
